%=====================================================================
\section{One complete episode, step by step}
\label{sec:trace}

This section follows the trained baseline model (\code{baseline\_RL/model/model\_best.pt}) through one episode on a real instance, with every number taken from the code (\code{report/trace\_episode.py}, which repeats \code{rollout} and records everything). The instance is \textbf{EV-C101 with its first 8 cars}: $T=1236$ minutes, 12 stations with 2 guns each, 5 vans, typical cost per car $\sigma=916.8$ Rs. The episode is one greedy rollout (map symmetry 0, the most likely action at every step).

\subsection{The cars}

\begin{table}[H]
\centering
\caption{The 8 cars: segment, income, request minute $r_i$, energy $E_i$ (kWh), range $R_i$ (km), reachable stations, queue tolerance $W_i$ (min), money sensitivity $\mu_i$, values of time $c^Q_i$, $c^C_i$, $c^V_i$ (Rs/min), charging minutes at a gun and by a van.}
\footnotesize\begin{adjustbox}{max width=\linewidth}\begin{tabular}{@{}rlrrrrlrrrrrrr@{}}\toprule
EV & segment & $\theta$ & $r_i$ & $E_i$ & $R_i$ & reachable FCS & $W_i$ & $\mu_i$ & $c^Q_i$ & $c^C_i$ & $c^V_i$ & $\tau^F_i$ & $\tau^M_i$ \\\midrule
1 & commercial & 0.40 & 912 & 15.5 & 27.7 & 1,2,3,4,5,6,7,8,9,10,11,12 & 30 & 1.10 & 4.22 & 3.31 & 3.31 & 42 & 31 \\
2 & commercial & 0.27 & 825 & 27.8 & 33.6 & 1,2,3,4,5,6,7,8,9,10,11,12 & 33 & 1.23 & 3.51 & 2.74 & 2.74 & 76 & 56 \\
3 & personal & 0.48 & 65 & 15.5 & 0.2 & none (stranded) & 38 & 1.02 & 13.37 & 5.35 & 2.67 & 42 & 31 \\
4 & commercial & 0.30 & 727 & 15.5 & 13.8 & 2,3,4,5,6,8,9,10,11,12 & 33 & 1.20 & 3.65 & 2.85 & 2.85 & 42 & 31 \\
5 & commercial & 0.39 & 15 & 15.5 & 16.4 & 1,2,3,4,5,6,7,8,9,10,11,12 & 30 & 1.11 & 4.15 & 3.26 & 3.26 & 42 & 31 \\
6 & commercial & 0.44 & 621 & 21.6 & 18.4 & 1,2,3,4,5,6,7,8,9,10,11,12 & 29 & 1.06 & 4.43 & 3.49 & 3.49 & 59 & 43 \\
7 & personal & 0.82 & 170 & 21.6 & 29.4 & 1,2,3,4,5,6,7,8,9,10,11,12 & 23 & 0.68 & 20.93 & 8.37 & 4.19 & 59 & 43 \\
8 & commercial & 0.17 & 255 & 21.6 & 10.9 & 2,3,4,5,6,8,10,11,12 & 36 & 1.33 & 2.93 & 2.26 & 2.26 & 59 & 43 \\
\bottomrule\end{tabular}\end{adjustbox}

\end{table}

Car 3 is stranded (range 0.17 km, no station reachable). Car 7 belongs to a rich private owner ($\theta=0.82$): their station minute is worth Rs 8.37, a van minute only Rs 4.19. The commercial cars have $c^C=c^V$ (downtime rate).

\subsection{The decisions}

\begin{table}[H]
\centering
\caption{The 8 steps. ``open'' = cars with at least one allowed resource; the three largest car probabilities $p(i)$; the chosen car and resource with their probabilities; start, finish, cost increase and the running total.}
\footnotesize\begin{adjustbox}{max width=\linewidth}\begin{tabular}{@{}rrlllrrrr@{}}\toprule
step & open & top EV probabilities & chosen EV ($p$) & chosen resource ($p$) & $S_i$ & $C_i$ & $\Delta Z$ (Rs) & cost so far \\\midrule
1 & 8 & 3: 1.00, 5: 0.00, 7: 0.00 & EV3 (1.00) & van1 (1.00) & 65.0 & 96.0 & 2287.4 & 2287.4 \\
2 & 7 & 5: 0.96, 7: 0.04, 8: 0.00 & EV5 (0.96) & FCS6-g1 (0.69) & 17.2 & 59.5 & 540.6 & 2828.0 \\
3 & 6 & 7: 1.00, 8: 0.00, 6: 0.00 & EV7 (1.00) & van1 (0.95) & 170.0 & 213.3 & 813.7 & 3641.7 \\
4 & 5 & 8: 1.00, 6: 0.00, 2: 0.00 & EV8 (1.00) & FCS12-g1 (0.53) & 256.5 & 315.5 & 758.9 & 4400.6 \\
5 & 4 & 6: 0.88, 2: 0.11, 4: 0.00 & EV6 (0.88) & FCS6-g2 (0.38) & 624.2 & 683.2 & 752.0 & 5152.6 \\
6 & 3 & 2: 0.97, 4: 0.03, 1: 0.00 & EV2 (0.97) & FCS6-g1 (0.50) & 828.7 & 904.4 & 973.9 & 6126.5 \\
7 & 2 & 4: 0.95, 1: 0.05 & EV4 (0.95) & FCS12-g2 (0.76) & 732.5 & 774.8 & 557.3 & 6683.7 \\
8 & 1 & 1: 1.00 & EV1 (1.00) & FCS12-g1 (0.40) & 920.5 & 962.8 & 568.2 & 7251.9 \\
\bottomrule\end{tabular}\end{adjustbox}

\end{table}

\paragraph{Step 1: the stranded car first.} Car 3 has probability 1.00. The network learned (first from the greedy rule, which puts stranded cars first, and then by reinforcement) that a car only a van can reach must be placed while a van is free. For car 3 every gun is masked (``out of range''); vans 2--5 are masked by the symmetry rule (identical unused vans), so only van 1 is offered and $p(\text{van1}\mid3)=1$. The van needs 16.1 minutes to reach the car, which requested at minute 65, so $S_3=65$, no wait, $C_3=65+31.02=96.02$. The cost increase:
\[
\begin{aligned}
\Delta Z(3,\text{van1})&=\underbrace{\mu_3p^\Mcs E_3}_{1.017\times25\times15.51=394.30}+\underbrace{\pi^\Mcs E_3/\eta}_{10.7\times15.51/0.9=184.40}+\underbrace{c^V_3\tau^\Mcs_3}_{2.674\times31.02=82.95}+\underbrace{c^V_3q^\Mcs_3}_{0}\\
&\quad+\underbrace{c^\Mcs(d_{03}+d_{30}-d_{00})}_{15\times(4.192+4.192)=125.77}+\underbrace{c^\Fcs}_{1500}=\textbf{2287.42 Rs}.
\end{aligned}
\]
The state after the step: van 1 at node 3, free at minute 96.02, battery $150-(15.51/0.9+0.25\times4.192)=131.72$ kWh, 1 seat used. Note that the van's way home is in the cost ($d_{30}$) but not yet taken from the battery: it was only checked (need $=15.51/0.9+0.25\times(4.192+4.192)=19.33\le150$).

\paragraph{Step 3: a rich private owner prefers the van.} Table~\ref{tab:step3} lists every resource the policy could choose for car 7.

\begin{table}[H]
\centering
\caption{Step 3, car 7: every resource, why it is masked or allowed, its start, wait, finish, cost increase and the policy's probability. Gun 2 of a station is listed only when it is not masked by the symmetry rule (at FCS6, gun 1 is busy until minute 59.5, so the two guns differ and both are offered).}
\label{tab:step3}
\footnotesize\begin{adjustbox}{max width=\linewidth}\begin{tabular}{@{}llrrrrr@{}}\toprule
resource & mask & start $S$ & wait & finish $C$ & $\Delta Z$ (Rs) & $p(k\mid i)$ \\\midrule
FCS1-g1 & ok & 230.2 & 0.0 & 289.2 & 1701.6 & 0.000 \\
FCS2-g1 & ok & 194.7 & 0.0 & 253.7 & 1231.1 & 0.000 \\
FCS3-g1 & ok & 196.2 & 0.0 & 255.2 & 1250.1 & 0.000 \\
FCS4-g1 & ok & 202.6 & 0.0 & 261.6 & 1336.1 & 0.000 \\
FCS5-g1 & ok & 190.0 & 0.0 & 249.0 & 1168.6 & 0.000 \\
FCS6-g1 & ok & 172.8 & 0.0 & 231.8 & 940.8 & 0.002 \\
FCS6-g2 & ok & 172.8 & 0.0 & 231.8 & 940.8 & 0.028 \\
FCS7-g1 & ok & 226.6 & 0.0 & 285.6 & 1654.3 & 0.000 \\
FCS8-g1 & ok & 197.9 & 0.0 & 256.9 & 1272.7 & 0.000 \\
FCS9-g1 & ok & 213.0 & 0.0 & 272.0 & 1473.6 & 0.000 \\
FCS10-g1 & ok & 204.9 & 0.0 & 263.9 & 1365.7 & 0.000 \\
FCS11-g1 & ok & 206.2 & 0.0 & 265.2 & 1384.0 & 0.000 \\
FCS12-g1 & ok & 174.2 & 0.0 & 233.2 & 958.2 & 0.019 \\
van1 & ok & 170.0 & 0.0 & 213.3 & 813.7 & 0.950 \\
van2 & ok & 170.0 & 0.0 & 213.3 & 2431.2 & 0.000 \\
van3 & same as unused van & -- & -- & -- & -- & 0.000 \\
van4 & same as unused van & -- & -- & -- & -- & 0.000 \\
van5 & same as unused van & -- & -- & -- & -- & 0.000 \\
\bottomrule\end{tabular}\end{adjustbox}

\end{table}

The cheapest gun is FCS6 (940.81 Rs), the van already on the road is cheaper (813.68 Rs):
\[
\begin{aligned}
\Delta Z(7,\text{FCS6})&=\underbrace{2.02}_{\mu c^{EV}d}+\underbrace{220.85}_{\mu p^\Fcs E}+\underbrace{188.18}_{\pi^\Fcs E}+\underbrace{35.79}_{c^Dt}+\underbrace{0}_{c^Qq}+\underbrace{493.97}_{c^C\tau^\Fcs}=940.81,\\
\Delta Z(7,\text{van1})&=\underbrace{368.09}_{\mu p^\Mcs E}+\underbrace{257.16}_{\pi^\Mcs E/\eta}+\underbrace{181.12}_{c^V\tau^\Mcs}+\underbrace{0}_{c^Vq^\Mcs}+\underbrace{7.31}_{c^\Mcs(d_{37}+d_{70}-d_{30})}=813.68 .
\end{aligned}
\]
At the station the owner stands by the car for 59 minutes at Rs 8.37 a minute (494 Rs); with the van the car is charged in 43 minutes while the owner is free, at Rs 4.19 a minute (181 Rs). This \emph{driver-free benefit} (313 Rs) outweighs the dearer van tariff (+147 Rs felt) and the dearer van energy (+69 Rs). The van's detour is tiny (0.49 km, 7.31 Rs) because car 7 lies almost on van 1's way home. A \emph{new} van (van 2) would cost 2\,431.17 Rs because of the 1\,500 Rs dispatch cost, which is why the policy gives it probability 0. The policy chooses van 1 with probability 0.95. The battery check: need $=21.63/0.9+0.25\times(d_{37}+d_{70})=25.20\le131.72$.

\paragraph{Steps 2 and 4--8: stations.} The commercial cars go to stations close to them (FCS6 and FCS12). No car ever waits: the requests are spread over the day, and when a gun is busy (FCS6 gun 1 after car 5, FCS12 gun 1 after car 8) the later car arrives after it is free. The order of decisions does not follow the request times (car 4, requested at minute 727, is placed after car 2, requested at 825): the policy may place cars in any order, and the simulator still gives each its earliest start. In steps 7 and 8 the policy did not take the cheapest option (557.3 Rs against 539.1 Rs, and 568.2 against 541.6): one greedy rollout is not always optimal, which is why solving tries 520 plans.

\subsection{The check}

\begin{table}[H]
\centering
\caption{The six parts of $Z$ recomputed by the independent checker (\code{check\_plan\_f4}), against the sum of the eight cost increases of the simulator.}
\small\begin{tabular}{@{}lr@{}}\toprule term & checker (Rs) \\\midrule
$Z_1$ & 29.54 \\
$Z_2$ & 133.09 \\
$Z_3$ & 1500.00 \\
$Z_4$ & 1464.24 \\
$Z_5$ & 2840.10 \\
$Z_6$ & 1284.97 \\
\midrule $Z$ (checker) & 7251.93 \\
 sum of $\Delta Z$ (simulator) & 7251.93 \\
\bottomrule\end{tabular}

\end{table}

\begin{itemize}
  \item The sum of the eight $\Delta Z$ equals the checker's $Z$ to the last paisa (7\,251.93 Rs), as Proposition~\ref{prop:sum} says. The checker found no broken rule.
  \item The van-km telescoping: van 1's route is depot $\to$ car 3 $\to$ car 7 $\to$ depot. The van-km parts of the two steps were $8.385$ km (step 1) and $0.488$ km (step 3); their sum, $8.872$ km, is exactly the length of the route, and $15\times8.872=133.09=Z_2$.
  \item $Z_3=1500$: one van used, charged once (at step 1).
\end{itemize}

\subsection{Against the MILP and the full RL solve}

The same instance solved by the MILP on the laptop: \textbf{optimal, $Z=7\,207.21$ Rs, in 0.39 s}. Its plan uses the same van route (cars 3 and 7) and puts the other six cars at FCS6 and FCS12. The traced greedy rollout is 0.62\,\% dearer (7\,251.93): its $Z_1$ and $Z_6$ are higher because of the two non-cheapest choices of steps 7 and 8; $Z_2$--$Z_5$ are identical. The full RL solve of Section~\ref{sec:solve} (520 episodes) finds a plan of \textbf{7\,208.84 Rs, 0.02\,\% above the optimum}, again with the same van route. On this small instance the greedy rule happens to find the optimum too.

\inshort{On a real instance every decision can be followed and every rupee accounted for: the stranded car gets the van first (paying the dispatch cost), the rich private owner joins that van because the driver-free benefit outweighs the van premium, the taxis go to near stations, and the eight cost increases add up exactly to the checker's $Z$. The best of 520 tries is within 0.02\,\% of the MILP optimum.}
